\section{Scope, Limits, and Price Corroboration}
\label{sec:price_scope_conclusion}

Section~\ref{sec:screening_forensics} carries the paper's
central enforcement-design result. Price evidence enters this
section in a narrower, but now substantive, role. The screen is
not designed to estimate cartel markups, damages, or welfare
losses. Overcharge measurement has its own evidentiary tradition
\citep{bryant1991price,connor2007price}, distinct from screening
\citep{marshall2012economics,harrington2008detecting}. What the
price record can do, however, is reveal the empirical signature
of the cover-bidding mechanism that motivates the screen. The
naive FL-price association is positive (+\valBetaBroad); under
overlap-cell ATT weighting it reverses to negative
(\valScopeOverlapCoef, $p\valScopeOverlapPlt$).
This sign reversal is not a specification artifact. It is the
predicted decomposition into a positive selection component
(cartels concentrate in cells with structurally higher
underlying prices) and a negative within-cell mechanism component
(cover bidders depress the observed winner price relative to the
reference price by manufacturing the appearance of competition).
Both components are documented empirically in this section.

This distinction matters for the paper's legal-economic
position. The decomposition does not deliver a damages estimate
because the observed price under cover bidding is a depressed
artifact of the theater, not an undistorted measurement of the
cartel's true overcharge. The decomposition does, however,
establish that the FL screen identifies environments where the
underlying cover-bidding mechanism is operative. The price
record therefore disciplines the legal-economic interpretation
in two directions: it rules out a naive damages reading (the
positive baseline coefficient is selection, not effect), and it
rules out the dismissal of price evidence as uninformative (the
within-cell sign reversal is consistent with the mechanism the
screen targets).

The section reports the broad descriptive price imprint
(\S\ref{sec:price_imprint_submission}), the formal decomposition
into selection and mechanism components
(\S\ref{sec:price_scope_submission}), the legal-economic
implications (\S\ref{sec:legal_economic_use_submission}), and
the adaptive deployment alternatives that follow from these
findings (\S\ref{sec:adaptive_deployment_submission}). The
institutional implication remains: the award-layer signal is a
trigger for richer evidence collection, not a stand-alone measure
of price injury.

\subsection{Price Evidence as Scope}
\label{sec:price_imprint_submission}

The price regressions are reported to assess economic
non-neutrality, not to estimate price injury. If the award-layer
screen ranks environments where cartel-adjacent loser-side
participation is more likely, it is informative to ask whether
those environments also differ in negotiated unit prices. The
answer can only play a corroborative role. The estimates below
are descriptive associations, not cartel markups, causal effects
of frequent-loser presence, or inputs into a damages calculation.

In the broad item-level regressions, frequent-loser presence is
positively associated with log negotiated unit price. The
coefficient is $\valTabPricesColOne$ with item and year fixed
effects and $\valTabPricesColTwo$ after adding procuring-unit
fixed effects. The association is larger in pregão,
$\valTabPricesColThree$, than in convite, $\valTabPricesColFour$.
Across the main descriptive specifications, the implied price
range is $\valHeadlineRange$.

The modality split is useful as scope discipline, not mechanism
identification. If frequent losers were only mechanical quorum
fillers under the three-bidder rule, the price imprint should be
stronger in convite. Instead, the larger association appears in
pregão, the electronic-auction environment without the same formal
invitation structure. This contrast weighs against a simple
procedural-artifact interpretation, but it still does not identify
a causal procurement-price effect.

\begin{table}[htbp]
\centering
\caption{Descriptive Price Imprint}
\label{tab:price_imprint_submission}
\footnotesize
\begin{adjustbox}{max width=\textwidth,center}
\begin{threeparttable}
\begin{tabular}{p{5.4cm}ccc}
\toprule
Exercise & Estimate & Standard error & Use in the paper \\
\midrule
Broad sample, item and year FE & $\valTabPricesColOne$ &
$\valTabPricesSEOne$ & Economic non-neutrality \\
Broad sample, item, year, PBU FE & $\valTabPricesColTwo$ &
$\valTabPricesSETwo$ & Positive after buyer controls \\
Pregão subsample & $\valTabPricesColThree$ &
$\valTabPricesSEThree$ & Stronger in electronic auctions \\
Convite subsample & $\valTabPricesColFour$ &
$\valTabPricesSEFour$ & Smaller under quorum-prone procedure \\
\bottomrule
\end{tabular}
\begin{tablenotes}
\footnotesize
\item \textit{Notes:} Dependent variable is log negotiated unit
price. Estimates are descriptive associations in the BEC
item-level sample and are not interpreted as causal price
effects, cartel markups, overcharges, or damages estimates.
\end{tablenotes}
\end{threeparttable}
\end{adjustbox}
\end{table}

The price imprint therefore corroborates the economic relevance
of the priority set, but it does not carry the paper's central
claim. The central claim remains about allocating the next unit
of proof-producing effort. The next subsection uses the overlap
and segment results to make that boundary explicit.

\subsection{Decomposing the Sign Reversal}
\label{sec:price_scope_submission}

The sign reversal between the broad-sample coefficient
($+\valBetaBroad$) and the overlap-cell ATT coefficient
($\valScopeOverlapCoef$) is not a specification puzzle to
explain away or a boundary to be conservative around. The two
coefficients measure two different objects, and the gap between
them admits a clean economic decomposition. The broad sample
averages across cells; cartels with cover bidders are not
randomly distributed across cells but systematically concentrate
where the underlying price level is structurally higher. The
overlap-cell ATT removes that across-cell variation and isolates
the within-cell effect; within cells, the cover-bidding theater
mechanically depresses the observed winner price relative to the
reference price. The naive baseline mixes the two opposing
forces; the sign-reversed ATT spec isolates the within-cell
mechanism.

We document both components empirically. The decomposition
follows the standard separation between descriptive comparisons
and causal estimands \citep{imbens2015causal}; coefficient
stability and sensitivity logic reinforce the same caution
\citep{oster2019unobservable,cinelli2020making}.

\paragraph{Selection: cartels in high-price cells.}
Among non-treated items (items without frequent-loser presence),
mean log negotiated price rises monotonically across cell
FL-share quintiles: $\valSelTestQOne$ at Q1 (cells where ${\sim}1.4\%$
of items are FL-present) and $\valSelTestQFive$ at Q5 (cells where
${\sim}\valPctFortyone$ of items are FL-present). The gap is
$\Delta=\valSelTestDelta$ log-points across cells, a structural
ordering of which products and procurement environments cartels
choose to operate in. An item-level OLS regression of
non-treated log price on cell FL-share with all five cell
dimensions as marginal fixed effects yields a coefficient of
$\valSelTestCoefFullFE$ (SE $\valSelTestSEFullFE$, $N=\valSelTestN$).
The selection component is positive, large, and survives
partialling out each cell-dimension marginal effect.

\paragraph{Mechanism: cover-bidding theater within cell.}
Within overlap cells, frequent-loser presence depresses the
observed winner-to-reference price ratio by
$\valMechWinnerVsRef$ log-points (SE $\valMechWinnerVsRefSE$,
$p<10^{-30}$). The same regression with
$\log(\text{n\_firms})$ as control returns
$\valMechWinnerVsRefControlled$ (SE $\valMechWinnerVsRefControlledSE$,
not significant): the within-cell price effect operates entirely
through the bidder-count channel. Frequent-loser presence is
associated with $\valMechNFirms$ log-bidders ($\approx
\valMechNFirmsPct$ more bidders) within the same cell type
(SE $\valMechNFirmsSE$, $p<10^{-30}$); those additional bidders
mechanically compress the winning bid toward the reference
price. The cover-bidding theater is empirically detectable both
in the bidder-count inflation and in its downstream price
depression.

\paragraph{The bidder-count boundary.}
The two opposing forces exchange dominance at the
tender-density threshold. In sparse tenders (1--3 or 4--6
bidders), frequent-loser presence is associated with higher
prices ($\Delta=\valMechSparseBidderDelta$ and
$\valMechMidBidderDelta$ in the two quartiles respectively):
selection dominates because cover-bidding theater requires a
critical mass of bidders to be visible. In dense tenders
(11+~bidders), frequent-loser presence is associated with
\emph{lower} prices ($\Delta=\valMechDenseBidderDelta$):
the mechanism dominates because cover-bidder inflation is
operative. The 7--10 bidder bucket is the transition
($\Delta=\valMechMidHighBidderDelta$). The selection-mechanism
boundary is therefore observable at the tender-density
threshold itself.

Together, the two components yield the empirical signature of
cover-bidding theory: positive selection across cells, negative
mechanism within cell, and a sharp bidder-count transition
separating the two regimes.

Under overlap-cell ATT reweighting, the aggregate coefficient
reverses to $\valScopeOverlapCoef$; under propensity-score
trimming it becomes $\valScopePSCoef$. Only
$\valDroppedItemsShare$ of treated items lack a within-cell
counterfactual, so the reversal is not mainly a
dropped-observation result. It is a weighting result. The broad
price imprint and the overlap estimate put weight on different
parts of the procurement distribution, and those parts carry
different economic content: the broad spec averages
selection-into-high-price-cells together with the cover-bidding
theater that depresses observed prices, while the ATT spec
isolates the latter.

Table~\ref{tab:sign_reversal_decomposition_submission} summarizes
the decomposition.

\begin{table}[htbp]
\centering
\caption{Selection + Mechanism Decomposition of the FL-Price Sign Reversal}
\label{tab:sign_reversal_decomposition_submission}
\footnotesize
\begin{adjustbox}{max width=\textwidth,center}
\begin{threeparttable}
\begin{tabular}{p{4.8cm}cp{2.0cm}p{5.6cm}}
\toprule
Component & Estimate & SE / p & Interpretation \\
\midrule
Broad-sample baseline & $+\valBetaBroad$ & $p=\valBetaBroadP$ &
  Joint effect: selection $+$ mechanism \\
\midrule
\multicolumn{4}{l}{\textit{Across-cell selection (non-treated price on cell FL-share)}} \\
\hspace{6pt} Raw OLS & $\valSelTestCoefRaw$ & SE $\valSelTestSERaw$ &
  Cells where cartels concentrate are structurally high-priced \\
\hspace{6pt} + item-group + year FE & $\valSelTestCoefFE$ & SE $\valSelTestSEFE$ &
  Selection beyond marginal cell dimensions \\
\hspace{6pt} + all five marginal cell-dim FE & $\valSelTestCoefFullFE$ &
  SE $\valSelTestSEFullFE$ &
  Selection coefficient after partialling out each cell dim \\
\midrule
\multicolumn{4}{l}{\textit{Within-cell mechanism (winner-to-reference ratio)}} \\
\hspace{6pt} Within cell, no controls & $\valMechWinnerVsRef$ &
  SE $\valMechWinnerVsRefSE$, $p<10^{-30}$ &
  Cover bidders compress winner-to-reference ratio \\
\hspace{6pt} + $\log(\text{n\_firms})$ control & $\valMechWinnerVsRefControlled$ &
  SE $\valMechWinnerVsRefControlledSE$, n.s. &
  Effect operates through bidder-count channel \\
\hspace{6pt} M1: log-bidder inflation & $\valMechNFirms$ &
  SE $\valMechNFirmsSE$, $p<10^{-30}$ &
  $\approx \valMechNFirmsPct$ more bidders in FL-present tenders \\
\midrule
\multicolumn{4}{l}{\textit{Sign-reversal boundary by bidder count}} \\
\hspace{6pt} 1--3 bidders (sparse) & $\valMechSparseBidderDelta$ & --- &
  Selection regime: FL $=$ higher price \\
\hspace{6pt} 7--10 bidders & $\valMechMidHighBidderDelta$ & --- &
  Transition \\
\hspace{6pt} 11+ bidders (dense) & $\valMechDenseBidderDelta$ & --- &
  Mechanism regime: FL $=$ lower price \\
\bottomrule
\end{tabular}
\begin{tablenotes}
\footnotesize
\item \textit{Notes:} Selection component: log negotiated price
among non-treated items only, regressed on cell FL-share
(across-cell variation). Mechanism component: log winner-to-
reference ratio regressed on FL presence with overlap-cell fixed
effects (within-cell variation). All standard errors clustered
by overlap cell. The bidder-count rows report
$\Delta(\text{FL-present} - \text{FL-absent})$ on winner-to-
reference within each bidder-count bucket. See Appendix~C and
AN-039, AN-040 in the supplementary site for the underlying
regressions.
\end{tablenotes}
\end{threeparttable}
\end{adjustbox}
\end{table}

The progression is monotone across the three specifications:
broad-sample $+\valBetaBroad$ (item, year, PBU fixed effects;
$p=\valBetaBroadP$); overlap-cell unweighted $+\valBetaUnwOverlap$
($p=\valBetaUnwOverlapP$); overlap-cell ATT $\valScopeOverlapCoef$
(within-cell ATT weighting, $p<10^{-9}$). The sign change occurs
between the unweighted overlap restriction and ATT reweighting,
not at the sample restriction itself. Within the overlap-ATT
specification the negative coefficient is not driven by a single
heterogeneous subset: it survives across both procurement
modalities (Convite $\valSignBetaConvATT$, $p<10^{-5}$; Pregão
$\valSignBetaPregATT$, $p<10^{-13}$), every PBU-size quartile, and
tender-value quartiles Q1 through Q3. The only systematic
positive cell within overlap-ATT is tender-value Q4, the segment
already discussed. The subgroup decomposition therefore
strengthens the scope reading rather than threatening it: the
negative reading appears in routine large-buyer environments,
while the high-value tail is where price evidence stops aligning
with the loser-side scope.

The segment decomposition makes the boundary sharper. In the
lowest three tender-value quartiles, the overlap estimates are
negative. In the highest tender-value quartile, the estimate
remains positive, $\valBetaQfour$ ($p=\valBetaQfourP$). That is
the segment where a cartel interpretation is most plausible ex
ante because fixed coordination costs are spread over larger
contracts. Even there, the result is scope evidence, not a
markup estimate. It tells the reader where the price imprint is
economically concentrated; it does not establish price injury.

\begin{table}[htbp]
\centering
\caption{Where Price Evidence Stops}
\label{tab:price_scope_submission}
\footnotesize
\begin{adjustbox}{max width=\textwidth,center}
\begin{threeparttable}
\begin{tabular}{p{5.2cm}ccp{4.8cm}}
\toprule
Exercise & Estimate & $p$ / SE & Interpretation \\
\midrule
Broad sample, item+year+PBU FE & $\valBetaBroad$ &
$p=\valBetaBroadP$ & Positive baseline before scope discipline \\
Overlap-cell unweighted & $\valBetaUnwOverlap$ &
$p=\valBetaUnwOverlapP$ & Attenuates under common-support restriction \\
Overlap-cell ATT & $\valScopeOverlapCoef$ &
$p\valScopeOverlapPlt$ & \textbf{Sign reversal under within-cell ATT weighting} \\
PS-trimmed ATT & $\valScopePSCoef$ &
$\valScopePSSE$ & Stronger negative estimate after trimming \\
Tender-value Q4, overlap ATT & $\valBetaQfour$ &
$p=\valBetaQfourP$ & Positive only in high-value segment \\
Direct-CADE overlap items & $\valBetaCADEDirect$ &
$p=\valBetaCADEDirectP$ & No separate price association for legal anchors \\
\bottomrule
\end{tabular}
\begin{tablenotes}
\footnotesize
\item \textit{Notes:} Dependent variable is log negotiated unit
price. The overlap rows restrict or reweight comparisons toward
cells with treated and untreated support. Estimates are used to
discipline the scope of price interpretation, not to estimate
causal effects, overcharges, markups, or damages.
\end{tablenotes}
\end{threeparttable}
\end{adjustbox}
\end{table}

The direct-CADE overlap estimate is also null,
$\valBetaCADEDirect$ ($p=\valBetaCADEDirectP$), which reinforces
the same interpretation. Price regressions do not recover the
legal object of the case, and they do not supply a damages base.
They show that the award-layer screen points to economically
non-neutral environments in parts of the distribution, while the
overlap reversal prevents price evidence from bearing more
weight than it can support.

\subsection{Implications for Legal-Economic Use}
\label{sec:legal_economic_use_submission}

The legal-economic use of the screen follows from these
boundaries. Price evidence may be consistent with economic
relevance, but it does not transform an award-layer statistic
into proof of agreement or price injury. The institutional use is staged:
award records rank where costly bid-level and documentary
evidence should be assembled; bid records, case documents, and
legal process determine whether the conduct supports an antitrust
case.

This distinction also clarifies portability. The architecture is
more portable than the estimates: BEC estimates should not be
exported as numerical external validity. What travels is the
information problem. Many
procurement systems observe participation and awards more cheaply
than they can recover, clean, and evaluate bid histories. In such
settings, a useful screen requires stable participant identities,
observable winners and losers, repeated procurement, sufficiently
comparable items, and a feasible path from the initial ranking to
richer forensic follow-up. Without that path to proof, a cheap
signal has little enforcement value.
Another agency should not import the BEC cutoff or performance
metrics mechanically; it should validate the ranking against its
own enforcement anchors and observability constraints.

\subsection{Adaptive Deployment}
\label{sec:adaptive_deployment_submission}

The screen should also be deployed as a moving queue rather than
as a static legal trigger. Any administrative rule can
generate strategic response once firms learn which observable
patterns attract attention. In this setting, firms could rotate
loser roles, engineer occasional wins, fragment identities, or
avoid a known participation threshold. That possibility does not
undermine the sequencing architecture, but it does rule out naive
use of the binary cutoff as a permanent enforcement boundary.
Strategic response weakens static cutoff use more than it
invalidates the two-stage architecture. It is a standard concern
for cartel enforcement and screening because firms may adapt to
observable enforcement rules once those rules are understood
\citep{green1984noncooperative,haltiwanger1991the}. Screening
rules face the same adaptation problem in enforcement settings
\citep{harrington2008detecting,sanchezgraells2019screening}.

The implementation implication is straightforward. Agencies
should use continuous ranks or capacity-constrained top-$k$
queues, refresh thresholds, monitor bunching below operational
cutoffs, and combine award-layer screening with bid-layer
forensics whenever bid files can be recovered. A firm that moves
to avoid one observable pattern may create another: discontinuous
participation, suspicious identity changes, or shifts in
bid-level behavior. The architecture is therefore strongest when
the award-layer queue is periodically updated and linked to
forensic follow-up rather than treated as a one-time blacklist.

This is also the right legal posture. The screen is not a rule of
liability and should not be used as one. It is a way to order
scarce investigative attention when proof is costly and the first
observable record is legally thin. Used adaptively, it gives an
agency a disciplined place to begin without pretending that the
beginning is the end of the evidentiary process.
